% !TEX TS-program = pdflatex
% !TEX encoding = UTF-8 Unicode

% This file is a template using the "beamer" package to create slides for a talk or presentation
% - Giving a talk on some subject.
% - The talk is between 15min and 45min long.
% - Style is ornate.

% MODIFIED by Jonathan Kew, 2008-07-06
% The header comments and encoding in this file were modified for inclusion with TeXworks.
% The content is otherwise unchanged from the original distributed with the beamer package.

\documentclass{beamer}
\setbeamercovered{invisible}



% Copyright 2004 by Till Tantau <tantau@users.sourceforge.net>.
%
% In principle, this file can be redistributed and/or modified under
% the terms of the GNU Public License, version 2.
%
% However, this file is supposed to be a template to be modified
% for your own needs. For this reason, if you use this file as a
% template and not specifically distribute it as part of a another
% package/program, I grant the extra permission to freely copy and
% modify this file as you see fit and even to delete this copyright
% notice. 


\mode<presentation>
{
  \usetheme{Warsaw}
  % or ...

  % or whatever (possibly just delete it)
}


\usepackage[english]{babel}
% or whatever

\usepackage[utf8]{inputenc}
\usepackage{qrcode}
\usepackage{hyperref}
\usepackage{ytableau}
\usepackage{longtable}
% or whatever

\usepackage{times}
\usepackage[T1]{fontenc}
\usepackage{tikz}
% Or whatever. Note that the encoding and the font should match. If T1
% does not look nice, try deleting the line with the fontenc.

%%% MATH RELATED

\title{MATH 180 Strategies of Problem Solving} 
\subtitle
{} % (optional)

\author[W.R. Casper] % (optional, use only with lots of authors)
{W.R. Casper}
% - Use the \inst{?} command only if the authors have different
%   affiliation.

\institute[California State University Fullerton] % (optional, but mostly needed)
{
  Department of Mathematics\\
  California State University Fullerton}
% - Use the \inst command only if there are several affiliations.
% - Keep it simple, no one is interested in your street address.

\subject{Talks}
% This is only inserted into the PDF information catalog. Can be left
% out. 



% If you have a file called "university-logo-filename.xxx", where xxx
% is a graphic format that can be processed by latex or pdflatex,
% resp., then you can add a logo as follows:

% \pgfdeclareimage[height=0.5cm]{university-logo}{university-logo-filename}
% \logo{\pgfuseimage{university-logo}}



% Delete this, if you do not want the table of contents to pop up at
% the beginning of each subsection:
\AtBeginSubsection[]
{
  \begin{frame}<beamer>{Outline}
    \tableofcontents[currentsection,currentsubsection]
  \end{frame}
}


% If you wish to uncover everything in a step-wise fashion, uncomment
% the following command: 

%\beamerdefaultoverlayspecification{<+->}


\begin{document}

\begin{frame}
  \titlepage
\end{frame}

\begin{frame}{Outline}
  \tableofcontents
  % You might wish to add the option [pausesections]
\end{frame}

% Since this a solution template for a generic talk, very little can
% be said about how it should be structured. However, the talk length
% of between 15min and 45min and the theme suggest that you stick to
% the following rules:  

% - Exactly two or three sections (other than the summary).
% - At *most* three subsections per section.
% - Talk about 30s to 2min per frame. So there should be between about
%   15 and 30 frames, all told.

\section{SoPS Lecture 20}

\begin{frame}{Review}
  \begin{block}{Last time}
    \begin{itemize}
\pause
      \item Relations, Functions
\pause
      \item Injective, Surjective, Bijective
\pause
      \item Cardinality, Hilbert's Hotel
    \end{itemize}
  \end{block}
\end{frame}

\subsection{Finite Sets}

\begin{frame}{Finite means what?}
\begin{block}{Question}
How should we define a set to be \textbf{finite}?
\end{block}
\end{frame}

\begin{frame}{Finite sets}
We have some sets we know are finite.\\
\pause
Define
$$\mathbb N_1 = \{1\},\quad\mathbb N_2 = \{1,2\},\quad\mathbb N_3 = \{1,2,3\}$$
\pause
In general, for any $k\in\mathbb N$ define
$$\mathbb N_k = \{1,2,3,\dots, k\}.$$
\pause
\begin{block}{Definition}
A set $A$ is called \textbf{finite} if it is empty or there exists a $k\in\mathbb N$ for which there is a bijection from $\mathbb N_k$ to $A$.
\end{block}
\end{frame}

\begin{frame}{Example}
\begin{block}{Problem}
Show that $\mathbb N_j$ is finite for all $j$.
\end{block}
\pause
Must show that there's a $k\in \mathbb N$ with a bijection from $\mathbb N_k$ to $\mathbb N_j$.\\
\pause
\begin{block}{Question}
What value of $k$ should we choose?
\end{block}
\pause
\begin{block}{Answer}
We should choose $k=j$.  That's the only value we should expect to have a bijection!
\end{block}
\end{frame}

\begin{frame}{Example}
\begin{block}{Problem}
Show that $\mathbb N_j$ is finite for all $j$.
\end{block}
Must show that there's a $k\in \mathbb N$ with a bijection from $\mathbb N_k$ to $\mathbb N_j$.\\
\begin{block}{Question}
What bijection is there from $\mathbb N_k$ to $\mathbb N_k$?
\end{block}
\pause
\begin{block}{Answer}
There's lots!  We just need to demonstrate that there's at least one.\\
\pause
The \textbf{identity function}
$$f:\mathbb N_k\rightarrow\mathbb N_k,\quad f(x) = x$$
in particular is a bijection.
\end{block}
\end{frame}

\begin{frame}{Example}
\begin{block}{Problem}
Show that $A = \{\heartsuit,\spadesuit,\diamondsuit,\clubsuit\}$ is finite.
\end{block}
Must show that there's a $k\in \mathbb N$ with a bijection from $\mathbb N_k$ to $A$.\\
\begin{block}{Question}
What $k$ should we choose?
\end{block}
\pause
\begin{block}{Answer}
We should choose $k=4$.
\end{block}
\end{frame}

\begin{frame}{Example}
\begin{block}{Problem}
Show that $A = \{\heartsuit,\spadesuit,\diamondsuit,\clubsuit\}$ is finite.
\end{block}
Must show that there's a $k\in \mathbb N$ with a bijection from $\mathbb N_k$ to $A$.\\
\begin{block}{Question}
What bijection from $\mathbb N_4$ to $A$ is there?
\end{block}
\pause
\begin{block}{Answer}
There's lots!  We just need to show that there's at least one.\\
\pause
For example, $f: \mathbb N_4\rightarrow A$ defined by
$$f(1) = \heartsuit,\quad f(2) = \spadesuit,\quad f(3) = \diamondsuit,\quad f(4) = \clubsuit$$
\end{block}
\end{frame}

\subsection{Pigeonhole Principle}
\begin{frame}{Starting Example}

\pause
\begin{block}{Question}
Why is there no bijection from $\mathbb N_2$ to $\mathbb N_1$?
\end{block}
\pause
\begin{block}{Solution}
\pause
Take a function $f: \mathbb N_2\rightarrow \mathbb N_1$.\\
\pause
$$f: \{1,2\}\rightarrow\{1\}$$
\pause
$$f(1) = 1\quad\text{and}\quad f(2)=1$$
\pause
Therefore $f$ fails to be injective.
\end{block}
\end{frame}

\begin{frame}{Second Example}
\begin{block}{Question}
Why is there no bijection from $\mathbb N_3$ to $\mathbb N_2$?
\end{block}
\pause
\begin{block}{Solution}
\pause
Take a function $f: \mathbb N_3\rightarrow \mathbb N_2$.\\
\pause
$$f: \{1,2,3\}\rightarrow\{1,2\}$$
\pause
If $f(1) = f(2)$, then not injective.\\
\pause
Otherwise either $f(1) = 1,f(2) = 2$ or $f(1) = 2,f(2) =1$.\\
\pause
In either case, $f(3)$ must repeat one of the values!\\
\pause
Therefore $f$ fails to be injective.
\end{block}
\end{frame}

\begin{frame}{Pigeonhole Principle}
\begin{block}{Theorem (Pigeonhole Principle)}
Let $j$ and $k$ be integers with $k > j$.
If $f: \mathbb N_k\rightarrow\mathbb N_j$, then $f$ is not injective.
\end{block}
\end{frame}

\begin{frame}{Application}
\begin{block}{Problem}
Show that at any party, two people in the party will always end up having the same number of acquaintances attending the party.
\end{block}
\pause
Fine print:
\begin{itemize}
\pause
\item we always consider someone to be their own acquaintance
\end{itemize}
\end{frame}

\begin{frame}{Application}
\begin{block}{Solution}
\pause
Label the guests at the party as $1,2,\dots,n$.\\
\pause
Define a function
$$f(k) = \text{number of people $k$ knows at the party}.$$
\pause
Consider two cases:
\begin{itemize}
\pause
\item nobody knows everyone at the party
\pause
\item somebody knows everyone at the party
\end{itemize}
\end{block}
\end{frame}

\begin{frame}{Application}
\pause
\begin{block}{Solution}
\pause
If nobody knows everyone at the party, then the values of $f$ are $1,2,\dots, n-1$.\\
\pause
Therefore $f$ is a function
$$f: \mathbb N_n\rightarrow\mathbb N_{n-1}.$$
\pause
The Pigeonhole Principle says that $f$ is not injective!\\
\pause
Therefore there are two different values $a$ and $b$ with $f(a) = f(b)$.\\
\pause
Guests $a$ and $b$ have the same number of acquaintances at the party!
\end{block}
\end{frame}

\begin{frame}{Application}
\begin{block}{Solution}
\pause
If somebody knows everyone at the party, then the values of $f$ are $2,3,\dots, n$.\\
\pause
Define $g(x) = f(x)-1$, which takes values $1,2,\dots, n-1$.\\
\pause
$$g: \mathbb N_n\rightarrow\mathbb N_{n-1}.$$
\pause
The Pigeonhole Principle says that $g$ is not injective!\\
\pause
Therefore there are two different values $a$ and $b$ with $g(a) = g(b)$.\\
\pause
This means $f(a)-1=f(b)-1$, so $f(a) = f(b)$
\pause
Guests $a$ and $b$ have the same number of acquaintances at the party!
\end{block}
\end{frame}

\begin{frame}{Application}
\begin{block}{Problem}
Let $p=8675309$.  Prove that there is a number $n$ of the form 
$$n = 111111\dots 1$$
which is divisible by $p$.
\end{block}
\pause
Observation:
\begin{itemize}
\pause
\item the number $p$ is prime, which might explain why Jenny didn't want to change it.
\end{itemize}
\end{frame}

\begin{frame}{Application}
\pause
\begin{block}{Solution}
Define $n_1 = 1$, $n_2 = 11$, $n_3 = 111$, and so on.\\
For each $k$, define $r_k$ to be the remainder of $n_k$ divided by $p$.
\pause
$$B = \{r_1,r_2,\dots, r_{p+1}\}\quad\text{has at most $p$ elements}$$
\pause
Pigeonhole principle: there's $j < k$ with $r_j = r_k$.\\
\pause
$$p\ \text{divides}\ n_k-n_j = \overbrace{111\dots 1}^{\text{$(k-j)$ many}}\overbrace{00\dots 0}^{\text{$j$ many}} = n_{k-j}\cdot10^{j}$$
\pause
Therefore $p$ divides $n_{k-j} = \overbrace{111\dots 1}^{\text{$(k-j)$ many}}$
\end{block}
\end{frame}

\begin{frame}{Natural numbers are infinite}
\pause
\begin{block}{Theorem}
The set $\mathbb N$ of natural numbers are infinite.
\end{block}
\pause
\begin{block}{Proof}
To see this, suppose that $\mathbb N$ was finite.\\
\pause
Then there would be an integer $k\in\mathbb N$ and a bijection between $\mathbb N$ and $\mathbb N_k$ for some $k\in\mathbb N$.
\pause
$$f: \mathbb N\rightarrow\mathbb N_k\quad\text{bijection}$$
\pause
Use $f$ to define a new function
$$g: \mathbb N_{k+1}\rightarrow\mathbb N_k,\quad g(k) = f(k).$$
\end{block}
\end{frame}

\begin{frame}{Natural numbers are infinite}
\pause
\begin{block}{Theorem}
The set $\mathbb N$ of natural numbers are infinite.
\end{block}
\pause
\begin{block}{Proof}
\pause
If $g(x) = g(y)$, then $f(x) = f(y)$.\\
\pause
Since $f$ is bijective, it's injective, so $x=y$.\\
\pause
This shows $g$ is injective.\\
\pause
Waaaaait a minute!  This violates the Pigeonhole Principle!\\
\pause
Our original assumption that $\mathbb N$ was finite must be false.
\end{block}
\end{frame}


\subsection{Wrapping up}

\begin{frame}{Final thoughts}
  \begin{block}{Reflection}
    \begin{itemize}
\pause
      \item  Give the formal definition of a finite set.
      \item  Explain why there's a bijection between $\mathbb N_j$ and $\mathbb N_k$ if and only if $j=k$.
    \end{itemize}
  \end{block}
\begin{center}
Submit your responses by \textbf{Tonight!}
\end{center}
\end{frame}


\end{document}


